\documentclass[11pt]{article}
\usepackage[T1]{fontenc}
\usepackage[utf8]{inputenc}
\usepackage{lmodern}
\usepackage[a4paper,margin=1in]{geometry}
\usepackage{microtype}
\usepackage{amsmath,amssymb,amsthm,mathtools}
\usepackage{booktabs,array}
\usepackage[dvipsnames]{xcolor}
\usepackage[colorlinks=true,linkcolor=MidnightBlue,citecolor=MidnightBlue,urlcolor=MidnightBlue]{hyperref}
\usepackage{enumitem}
\usepackage{setspace}
\usepackage{fancyhdr}
\usepackage{titlesec}
\setlength{\headheight}{14pt}
\pagestyle{fancy}
\fancyhf{}
\fancyhead[L]{Sibling-interface descriptors v2 for $\mathcal{ALCI}_{RCC5}$}
\fancyhead[R]{May 2026}
\fancyfoot[C]{\thepage}
\titleformat{\section}{\large\bfseries\color{MidnightBlue}}{\thesection}{0.6em}{}
\titleformat{\subsection}{\normalsize\bfseries\color{MidnightBlue}}{\thesubsection}{0.5em}{}
\setlist[itemize]{topsep=4pt,itemsep=2pt,leftmargin=1.5em}
\setlist[enumerate]{topsep=4pt,itemsep=2pt,leftmargin=1.8em}
\onehalfspacing

\newtheorem{theorem}{Theorem}[section]
\newtheorem{proposition}[theorem]{Proposition}
\newtheorem{lemma}[theorem]{Lemma}
\newtheorem{corollary}[theorem]{Corollary}
\theoremstyle{definition}
\newtheorem{definition}[theorem]{Definition}
\newtheorem{construction}[theorem]{Construction}
\theoremstyle{remark}
\newtheorem{remark}[theorem]{Remark}

\newcommand{\DR}{\mathsf{DR}}
\newcommand{\PO}{\mathsf{PO}}
\newcommand{\PP}{\mathsf{PP}}
\newcommand{\PPI}{\mathsf{PPI}}
\newcommand{\EQ}{\mathsf{EQ}}
\newcommand{\RCC}{\{\DR,\PO,\PP,\PPI\}}
\newcommand{\cl}{\operatorname{cl}}
\newcommand{\Tp}{\operatorname{Tp}}
\newcommand{\State}{\operatorname{State}}
\newcommand{\Child}{\operatorname{Child}}
\newcommand{\Iface}{\operatorname{Iface}}
\newcommand{\Wit}{\operatorname{Wit}}
\newcommand{\Par}{\operatorname{Par}}
\newcommand{\Need}{\operatorname{Need}}
\newcommand{\st}{\operatorname{st}}
\newcommand{\stat}{\operatorname{stat}}
\newcommand{\comp}{\operatorname{comp}}
\newcommand{\inv}{\operatorname{inv}}
\newcommand{\orig}{\operatorname{orig}}
\newcommand{\Core}{\mathsf{Core}}
\newcommand{\Out}{\mathsf{Out}}
\newcommand{\Front}{\mathsf{Front}}
\newcommand{\Dom}{\mathsf{Dom}}
\newcommand{\OpenDom}{\mathbb{D}}
\newcommand{\md}{\operatorname{md}}
\newcommand{\Can}{\operatorname{Can}}
\newcommand{\gen}{\mathsf{gen}}
\newcommand{\reach}{\rightsquigarrow}
\newcommand{\ALCIRCC}[1]{\mathcal{ALCI}_{\mathsf{RCC#1}}}

\begin{document}

\begin{center}
{\LARGE\bfseries Finite Sibling-Interface Descriptors for $\ALCIRCC{5}$ (v2, Repaired)\par}
\vspace{0.5em}
{\large Reinterpreted split-tree presentation, ambient root removed,\\
 PP/PPI eventualities by reachability discharge}\par
\vspace{0.75em}
{\normalsize Claude (Anthropic, Opus 4.7), prompted by Michael Wessel}\par
\vspace{0.45em}
{\small May 2026}
\end{center}

\begin{center}
\fbox{\parbox{0.92\textwidth}{%
\textbf{Disclaimer.}\;
This document was generated by Claude (Anthropic), prompted by Michael
Wessel. The results have not been independently verified. We invite
scrutiny and corrections.}}
\end{center}

\begin{abstract}
This is a repaired version of GPT-5.4's sibling-interface
note~\cite{SiblingV1} (April 2026), which proves the soundness
direction of the split-forest decidability programme for $\ALCIRCC{5}$.
GPT-5.5's review~\cite{GPT55Review} (May 2026) found that the v1 paper's
soundness chain (Thms 1.17--1.19) is intact, but identified two
structural defects elsewhere in the document: (i) the split-tree
presentation~\cite[Def 1.5]{SiblingV1} requires an
\emph{ambient root} above the evaluation point, which fails for
formulas forcing infinite $\PP$-chains
(e.g.\ $C_\uparrow := \exists\PP.\top \sqcap \forall\PP.\exists\PP.\top$);
and (ii) the proof sketch of Thm 1.20~\cite[Thm 1.20 sketch]{SiblingV1}
extracts ``least locally coherent support-closed tables'' without
constructing them and conflates witnesses across distinct occurrences
when discharging the C5 triple-coherence axiom.

This v2 paper applies three small but conceptually significant
repairs while preserving every soundness statement of v1 verbatim.
First, the split-tree presentation is reinterpreted as a
\emph{generated} witness-cover presentation (parent function may
self-loop at the root, so $C_\uparrow$ is represented by a
single-state cycle on the parent function). Second, $\PP/\PPI$
eventualities are discharged by \emph{finite-graph reachability}
on the quotient's parent graph, replacing the witness-menu
treatment of v1 line~367 which left transitive eventualities
under-specified. Third, the Thm 1.20 proof sketch is dropped: the
model-to-quotient direction is delegated to the explicit
construction in the companion v2 completeness-extraction
paper~\cite{CompletenessV2}, which builds the support-closed
mosaics that the v1 sketch only asserted.

The soundness chain $\text{valid finite quotient} \to
\text{disjunctive network} \to \text{arc-consistency} \to
\text{K\"onig} \to \text{canonical refinements} \to
\text{weak-}\EQ\text{ model} \to \text{strong-}\EQ\text{ model}$
(v1 Thms 1.17--1.19) is unchanged. Combined with v2 completeness
extraction, this yields a clean decidability theorem for
concept satisfiability in $\ALCIRCC{5}$ under strong-$\EQ$
semantics, without any $\omega$-acceptance machinery.
\end{abstract}

\tableofcontents

\section{What this v2 paper is and is not}
\label{sec:what-this-is}

\subsection{Pointer to v1 for unchanged material}

Every definition and theorem of GPT-5.4's v1 paper~\cite{SiblingV1}
is retained verbatim except where this note explicitly replaces or
reinterprets it. Specifically, the following v1 material is
unchanged and is cited rather than reproduced:

\begin{itemize}
\item v1 \S2 (RCC5 conventions, types, modal depth, type-safe open
  labels, per-relation $\Need_R$ families).
\item v1 \S4 (the corrected status partition: $\Core$, $\Out$,
  $\Front$, partition lemma).
\item v1 \S5 (rigid $\DR$ behaviour: uniform-$\DR$ lemma and the
  open cross-edge lemma).
\item v1 \S6 (status evolution; rank-0 colors; endpoint summaries;
  rank-$k$ descriptors; local coherence axioms (C1)--(C5)).
\item v1 \S7 (rank-$(k+1)$ states, except for the
  $\Wit_k$ component, which is reinterpreted in
  Section~\ref{sec:eventualities} below; the finite-index
  Theorem 1.10 is unchanged).
\item v1 \S8 (valid finite rank-$d$ quotient, unfolding $U(Q)$,
  disjunctive network of a finite prefix).
\item v1 \S9 (finite-prefix arc-consistency, Theorem 1.17).
\item v1 \S10 (prefix extension, canonical refinements,
  typed-$\EQ$ congruence, $\EQ$-synchrony).
\item v1 \S11 (weak-$\EQ$ realisation, Theorem 1.18; strong-$\EQ$
  realisation, Theorem 1.19).
\end{itemize}

\noindent\textbf{Replaced or reinterpreted.}\;
The following v1 items are altered:

\begin{itemize}
\item v1 Def 1.5 (weak-$\EQ$ split-tree presentation, line 158--172):
  the requirement of an \emph{ambient root}
  $\varrho_T$ is dropped; tree edges are reinterpreted as
  \emph{generated} witness-cover edges relative to a
  finite selected witness set, not as Hasse covers of an ambient
  $\PP$-order; the parent function may carry a self-loop at the root
  to represent infinite $\PP$-chains.
\item v1 Def 1.7 ($\Wit_k$ in the rank-$(k+1)$ state, line 367):
  $\PP/\PPI$ existentials are removed from the witness menu and
  discharged by reachability on the parent graph instead.
\item v1 Thm 1.20 (decidability sketch, line 681--691): the
  proof sketch is dropped and replaced by a citation to v2
  completeness extraction~\cite{CompletenessV2}, which carries out
  the model-to-quotient construction in full.
\end{itemize}

\subsection{What v2 does not change}

This note does not introduce any of the following:
\begin{itemize}
\item No B\"uchi automata, $\omega$-words, or parity acceptance:
  $\PP/\PPI$ eventualities are handled by ordinary finite-graph
  reachability, not by $\omega$-language theory.
\item No new model class: split-forest semantics, Hintikka types,
  RCC5 atomic relations, and strong-$\EQ$ models are exactly as
  in v1.
\item No new descriptor algebra: the local coherence axioms
  (C1)--(C5) are the v1 axioms, applied to the same rank-$k$
  descriptors. The witness-conflation defect that GPT-5.5
  identified in the v1 Thm~1.20 sketch is not a defect in the
  axioms themselves; it is a defect in how the v1 sketch
  \emph{extracted} the axiom-satisfying tables from a model. That
  extraction is carried out properly in v2 completeness
  extraction~\cite[\S 6, \S 9]{CompletenessV2}.
\end{itemize}

\subsection{Why a v2 paper rather than an in-place edit}

GPT-5.4 wrote the v1 paper in April 2026, before GPT-5.5's review
appeared. Editing it in place would mix attribution: the
reinterpreted definitions and the replaced Thm~1.20 sketch are
Claude's, not GPT-5.4's. v1 is therefore retained unchanged in
the repository, marked as superseded; v2 is this separate Claude
note. The arrangement parallels the
v1~\cite{CompletenessV1}/v2~\cite{CompletenessV2} split for the
companion completeness-extraction paper.

\section{The split-tree presentation, reinterpreted}
\label{sec:split-tree}

\subsection{The defect in v1 Def 1.5}

v1 Def 1.5 defines a weak-$\EQ$ split-tree presentation as a
rooted tree $T$ with an \emph{ambient root}~$\varrho_T$ above the
distinguished evaluation node~$r_\ast$. Tree edges are read as
``immediate $\PPI$ edges downward and hence $\PP$ upward,'' and
all indirect $\PP/\PPI$ edges are recovered from the ancestor
relation. This presupposes that every relevant model has a
$\PP$-maximal element above the evaluation point, in the same
sense as a Hasse-cover diagram of a well-founded order.

GPT-5.5's review~\cite{GPT55Review} identifies a counterexample:
\[
  C_\uparrow \;:=\;
  \exists\PP.\top \;\sqcap\; \forall\PP.\exists\PP.\top.
\]
Under the v1 orientation ($\PP$-witness = strict superpart), this
concept is satisfiable in any model whose domain contains a
strictly increasing $\PP$-chain
\[
  a_0 \;\PP\; a_1 \;\PP\; a_2 \;\PP\; \cdots
\]
with $C_\uparrow$ realised at $a_0$. No element of this domain is
$\PP$-maximal, so no ambient root~$\varrho_T$ exists. Yet the
concept is satisfiable. Therefore v1 Def~1.5 is too restrictive:
it does not admit a presentation for every satisfiable concept.

The defect is structural rather than cosmetic. It is what GPT-5.5
calls the \emph{rooted Hasse} assumption~\cite{GPT55Review,
GPT55Survival}.

\subsection{Reinterpretation: generated witness-cover presentations}

The fix preserves the tree-shaped containment backbone and the
$\PPI$-down/$\PP$-up reading, but drops the ambient-root
requirement and reinterprets tree edges as \emph{generated}
witness-cover edges relative to a finite selected witness set.
The construction is carried out in detail in v2 completeness
extraction~\cite[\S 4--\S 5]{CompletenessV2}; we summarise it here.

\begin{definition}[Generated weak-$\EQ$ split-tree presentation;
  v2 reinterpretation of v1 Def 1.5]\label{def:gen-split-tree}
A \emph{generated weak-$\EQ$ split-tree presentation} of a candidate
model of $C_0$ consists of:
\begin{enumerate}[label=(\roman*)]
  \item a tree $T$ of \emph{occurrences} with a distinguished
    \emph{evaluation node}~$r_\ast \in T$ (no ambient root is
    required);
  \item a partial parent function
    $\mathsf{par} : T \rightharpoonup T$, allowing
    $\mathsf{par}(u) = u$ at the root of an upward chain
    (\emph{parent self-loop});
  \item a map $\orig : T \to D$ onto an underlying generated
    submodel $D$ (defined below);
  \item a local type assignment $\lambda : T \to \Tp(C_0)$;
\end{enumerate}
such that:
\begin{itemize}
  \item tree edges are interpreted as immediate $\PPI$ edges
    downward and hence $\PP$ upward, \emph{relative to the
    finite witness set $\orig(T)$ used in the construction}; they
    are not claimed to be Hasse covers of any ambient $\PP$-order;
  \item if a node of $D$ must appear under multiple distinct
    parents, it may be represented by several occurrences in $T$
    (split copies);
  \item two occurrences are in relation $\EQ$ exactly when they
    are split copies of the same element of $D$ that satisfy the
    typed-$\EQ$ congruence of v1 Def~1.16.
\end{itemize}
A parent self-loop $\mathsf{par}(u) = u$ records that $u$
realises an upward $\PP$-existential whose witness has the same
rank-$d$ profile as $u$. The unfolding (Section~\ref{sec:unfolding})
expands such a self-loop into a fresh occurrence at every level.
\end{definition}

\begin{remark}[Why parent self-loops absorb $C_\uparrow$]
For $C_\uparrow$, every element of the witnessing chain
$a_0 \,\PP\, a_1 \,\PP\, \cdots$ has the same rank-$d$ profile
(local type $\{C_\uparrow, \exists\PP.\top, \forall\PP.\exists\PP.\top, \ldots\}$, plus the
$\Par_d$/$\Child_d$ multiplicities which are stationary along the
chain). The finite quotient therefore identifies all $a_i$ at a
single rank-$d$ state $\sigma$ and records the parent function as
$\mathsf{par}_Q(\sigma) = \sigma$. Unfolding this self-loop
produces the infinite chain $a_0 \PP a_1 \PP \cdots$ as required;
$C_\uparrow$'s satisfiability is thus represented by a quotient of
size~$1$ (one rank-$d$ state, one parent edge, a self-loop), with
no ambient root.
\end{remark}

\begin{remark}[Generated cover edges, in one paragraph]
\label{rem:generated-covers}
The phrase ``generated witness-cover edge''
\cite{GPT55Survival,CompletenessV2} means the following. Fix a
witness function $w$ that, for each occurrence $u$ and each
$\exists R.D \in \lambda(u)$, selects a model element witnessing
the existential. Let $A$ be the finite set of occurrences built
so far. A generated $\PPI$-cover edge $u \to v$ is added when
$\rho(\orig(u),\orig(v)) = \PPI$ and no other occurrence
$z \in A$ has $\rho(\orig(u),\orig(z)) = \PPI$ and
$\rho(\orig(z),\orig(v)) = \PPI$.  This is well-defined regardless
of whether the ambient $\PP$-order is dense or has Hasse covers:
the cover is taken \emph{relative to $A$}, not in the model. As
$A$ grows, new occurrences may slot between previously adjacent
ones; the construction proceeds inductively and is carried out
explicitly in~\cite[Construction 5.1]{CompletenessV2}.
\end{remark}

\subsection{Compatibility with v1 soundness theorems}

v1 Theorems 1.17--1.19 are statements about \emph{abstract valid
rank-$d$ quotients}: given a valid quotient $Q$ in the sense of
v1 Def~1.10, an unfolding $U(Q)$ is built, the disjunctive network
of every finite prefix has non-empty $\Need_R$-filtered domains,
canonical refinements yield a globally consistent atomic
labelling, and the resulting weak-$\EQ$ model is then quotiented
to a strong-$\EQ$ model satisfying $C_0$. None of these proofs
inspect whether the quotient came from a Hasse construction or a
generated-cover construction. They only use:
\begin{itemize}
  \item finiteness of $S$ (rank-$d$ states);
  \item finiteness of the disjunctive network alphabet;
  \item the local coherence axioms (C1)--(C5);
  \item the parent slot $\Par_d$ and child multiplicities
    $\Child_d$ as components of each state.
\end{itemize}
Allowing $\mathsf{par}_Q(\sigma) = \sigma$ does not break any of
these proofs: the unfolding step (v1 Def~1.11) constructs a fresh
occurrence at each predecessor level, and a self-loop simply
means that every level shares the same state. König's lemma in
v1 Lemma 1.13 still produces an infinite branch when one is
needed, and the canonical-refinements argument of v1 Lemma~1.16
treats every prefix uniformly.

\begin{lemma}[v1 soundness chain is preserved]\label{lem:soundness-preserved}
Replacing v1 Def~1.5 with the generated split-tree presentation of
Definition~\ref{def:gen-split-tree} (parent self-loops permitted,
no ambient root) leaves v1 Theorems 1.17, 1.18, and 1.19 valid
without modification.
\end{lemma}

\begin{proof}[Proof sketch]
Each of these theorems takes a valid finite rank-$d$ quotient $Q$
as input and constructs an RCC5 model. The proofs use the
finite-state quotient structure $(S, \mathsf{par}_Q, \Child_Q,
\Iface_Q)$, not Def~1.5 directly. The only place Def~1.5 enters is
when the unfolding $U(Q)$ is described as a tree of occurrences
with $\PPI$-down edges. With self-loops permitted, $U(Q)$ becomes
the regular unfolding of the parent graph: each
$\mathsf{par}_Q(\sigma) = \sigma$ self-loop expands into an
infinite branch of fresh occurrences with state $\sigma$. The
arc-consistency theorem (v1 Thm 1.12) iterates over finite
prefixes, all of which remain finite (each parent step adds one
occurrence). König's lemma (v1 Lemma 1.13) extracts a globally
consistent atomic labelling from compatibility on every finite
prefix. The canonical-refinements / typed-$\EQ$-congruence steps
(v1 Lemmas 1.15 and 1.16) operate on the resulting infinite
labelling. Strong-$\EQ$ realisation (v1 Thm 1.19) quotients by the
typed-$\EQ$ congruence. None of these steps need an ambient root.
\end{proof}

\subsection{Unfolding of a self-loop}
\label{sec:unfolding}

For completeness we record what the unfolding of a parent
self-loop looks like, since the v1 unfolding (Def~1.11) only
covered the rooted-tree case.

\begin{definition}[Unfolding $U(Q)$, v2]\label{def:UQ-v2}
Let $Q = (S, \mathsf{par}_Q, \Child_Q, \Iface_Q, \sigma_\ast)$ be
a valid finite rank-$d$ quotient with distinguished evaluation
state~$\sigma_\ast \in S$. The \emph{unfolding} $U(Q)$ is the
(possibly infinite) tree of occurrences defined coinductively:
\begin{enumerate}
\item create a root occurrence $r_\ast$ with $\st(r_\ast) =
  \sigma_\ast$;
\item for each occurrence $u$ with $\st(u) = \sigma$ and each
  child slot $\sigma' \in \Child_Q(\sigma)$ (with multiplicity
  $1$ or $+$), create one or more child occurrences with
  $\st = \sigma'$;
\item for each occurrence $u$ with $\st(u) = \sigma$, if
  $\mathsf{par}_Q(\sigma) = \sigma'$ for some $\sigma' \neq
  \sigma$, create a fresh parent occurrence $u^\uparrow$ with
  $\st(u^\uparrow) = \sigma'$ and set
  $\mathsf{par}_{U(Q)}(u) = u^\uparrow$;
\item if $\mathsf{par}_Q(\sigma) = \sigma$ (self-loop), create a
  fresh parent occurrence $u^\uparrow$ with $\st(u^\uparrow) =
  \sigma$ and set $\mathsf{par}_{U(Q)}(u) = u^\uparrow$. The same
  rule applies to $u^\uparrow$, producing an infinite chain of
  fresh occurrences all in state $\sigma$.
\end{enumerate}
\end{definition}

\begin{remark}[Self-loops are blocking, not equality]
A parent self-loop $\mathsf{par}_Q(\sigma) = \sigma$ records that
the rank-$d$ \emph{profile} repeats; it does \emph{not} identify
the two occurrences as $\EQ$-mates. The unfolding produces fresh
occurrences at every level. This is the distinction stressed
in~\cite{GPT55Survival}: blocking equivalence and typed-$\EQ$
congruence are different.  Only the latter justifies a quotient
of the unfolded model.
\end{remark}

\section{$\PP/\PPI$ eventualities by reachability discharge}
\label{sec:eventualities}

\subsection{The defect in v1 Def 1.7}

v1 Def~1.7 (rank-$(k+1)$ state, line~367) records, as part of
each state $\sigma$, a finite \emph{witness menu}~$\Wit_k(\sigma)$
recording, for every existential $\exists R.D \in \lambda(\sigma)$,
``one designated target color and target relation domain.'' The
patched cover-tree tableau note~\cite{CoverTreeTableauPatched} also
puts $\PP/\PPI$ existentials in the witness menu via $\mathsf{ReqUp}$
and $\mathsf{ReqDown}$ slots. By contrast, the v1 completeness
extraction~\cite{CompletenessV1} states (lines 399--401) that
$\exists\PP.D$ and $\exists\PPI.D$ are handled by the tree
structure and \emph{do not enter the witness menu}.

GPT-5.5~\cite{GPT55Review} flagged this inconsistency and pointed
out a real obstacle: a designated target color does not, by
itself, guarantee that the existential is realisable in the
unfolded model. For transitive roles, the witness for
$\exists\PP.D$ at $\sigma$ might lie arbitrarily far up the
parent chain --- in particular, beyond any finite parent slot.

\subsection{The fix: reachability over the parent graph}

The v2 completeness extraction~\cite[\S 11]{CompletenessV2}
discharges $\PP/\PPI$ eventualities by ordinary reachability over
the finite parent graph $(S, \mathsf{par}_Q)$.

\begin{definition}[Parent reachability;
  cf.\ \cite{CompletenessV2} Def 11.1]\label{def:parent-reach}
For $\sigma, \sigma' \in S$, write $\sigma \reach_{\PP} \sigma'$
when there exist $\sigma_0,\sigma_1,\ldots,\sigma_n$ in $S$ with
$\sigma_0 = \sigma$, $\sigma_n = \sigma'$, $n \ge 1$, and
$\mathsf{par}_Q(\sigma_i) = \sigma_{i+1}$ for $0 \le i < n$.
A self-loop $\mathsf{par}_Q(\sigma) = \sigma$ counts as $\sigma
\reach_{\PP} \sigma$. Symmetrically for $\reach_{\PPI}$ via the
inverse parent function (i.e., on child edges).
\end{definition}

\begin{definition}[Reachability discharge condition (v2 axiom)]
\label{def:reach-discharge}
A valid rank-$d$ quotient $Q$ must satisfy: for every $\sigma \in
S$ and every $\exists R.D \in \lambda_Q(\sigma)$ with $R \in
\{\PP, \PPI\}$, there exists $\sigma' \in S$ with
$\sigma \reach_R \sigma'$ and $D \in \lambda_Q(\sigma')$.
\end{definition}

\begin{remark}[v2 amendment to v1 Def 1.10]
v1 Def 1.10 (valid finite rank-$d$ quotient) is amended as
follows: replace the witness-menu condition for $\PP/\PPI$
existentials by Definition~\ref{def:reach-discharge}.
Existentials with $R \in \{\DR, \PO\}$ continue to be handled by
the witness menu $\Wit_k$ as in v1 Def~1.7. The witness menu
therefore retains its $\Need_R$-filtering side conditions for
$\DR$ and $\PO$; it simply no longer carries any $\mathsf{ReqUp}/
\mathsf{ReqDown}$ slots.
\end{remark}

\begin{remark}[This is finite-graph reachability, not B\"uchi]
Definition~\ref{def:parent-reach} is reachability in a finite
directed graph, decidable in $O(|S| + |E_Q|)$ per eventuality by
breadth-first search. There is no $\omega$-language theory and no
acceptance condition. A self-loop
$\mathsf{par}_Q(\sigma) = \sigma$ certifies that the parent chain
in the unfolding is infinite and remains in the strongly
connected component of $\sigma$; if some $\sigma'$ in that SCC
satisfies $D \in \lambda_Q(\sigma')$, the eventuality is
discharged. See~\cite[\S 11]{CompletenessV2} and
\cite[Remark 11.5]{CompletenessV2} for the comparison with
GPT-5.5's automata-theoretic
formulation~\cite[Def 6.4]{GPT55Generated}.
\end{remark}

\subsection{Compatibility with v1 soundness theorems}

The amendment to v1 Def~1.10 changes which quotients count as
valid; it does not change what valid quotients unfold to. v1
Thms 1.17--1.19 take a valid quotient as input and proceed
purely in terms of $S$, $\mathsf{par}_Q$, $\Child_Q$, and the
disjunctive networks of finite prefixes. None of those proofs
relies on $\Wit_k$ carrying $\mathsf{ReqUp}/\mathsf{ReqDown}$
slots. The arc-consistency theorem (v1 Thm 1.12) only needs
$\Need_R$-filtering for $R \in \{\DR, \PO\}$, which is the part of
$\Wit_k$ that v2 retains.

\begin{lemma}[Soundness chain unchanged]
\label{lem:soundness-chain-eventualities}
Replacing the v1 witness-menu treatment of $\PP/\PPI$
existentials by reachability discharge
(Definition~\ref{def:reach-discharge}) leaves v1 Theorems
1.17--1.19 valid without modification.
\end{lemma}

\begin{proof}[Proof sketch]
A valid rank-$d$ quotient under the v2 axioms is also valid under
the v1 axioms restricted to the part actually used by Thms~1.17--
1.19 (witness menus for $\DR/\PO$, plus all of (C1)--(C5),
$\Need_R$-filtering, and the parent/child slots).
$\PP/\PPI$ existentials are realised in the unfolding by a finite
parent-chain step: if $\sigma \reach_{\PP} \sigma'$ in $Q$ with $D
\in \lambda_Q(\sigma')$, the unfolding contains a finite branch
$u_0,\ldots,u_n$ with $\st(u_0) = \sigma$, $\st(u_n) = \sigma'$,
and $D \in \lambda(u_n)$. By v1 Lemma~1.4 (closure under modal
depth), the universal companions of $\exists\PP.D$ propagate
correctly along this chain, and the local-type Hintikka
conditions are preserved.
\end{proof}

\section{Decidability via the v2 chain}
\label{sec:decidability}

\subsection{Replacement for v1 Thm 1.20}

\begin{theorem}[Decidability of $\ALCIRCC{5}$ concept satisfiability;
  v2]\label{thm:decidability-v2}
Concept satisfiability in $\ALCIRCC{5}$ under strong-$\EQ$
semantics is decidable.
\end{theorem}

\begin{proof}
Combine three results.

\smallskip
\noindent\emph{Soundness} (valid quotient $\to$ model). By
v1~\cite[Thms 1.17, 1.18, 1.19]{SiblingV1}, every valid finite
rank-$d$ quotient $Q$ unfolds to an arc-consistent disjunctive
network whose König-extracted canonical refinement quotients to a
strong-$\EQ$ model of $C_0$.
Lemmas~\ref{lem:soundness-preserved}
and~\ref{lem:soundness-chain-eventualities} establish that the v2
reinterpretation of Def~1.5 (parent self-loops, generated covers)
and the v2 amendment of Def~1.10
($\PP/\PPI$-by-reachability-discharge) leave these theorems
valid.

\smallskip
\noindent\emph{Completeness} (model $\to$ valid quotient). By
v2 completeness extraction~\cite[Thm 12.1]{CompletenessV2}, every
pointed strong-$\EQ$ model $(\mathcal{I}, a_0) \models C_0$
extracts to a valid finite rank-$d$ quotient $Q$ whose
distinguished evaluation state contains $C_0$. The construction
uses witness-generated submodels~\cite[\S 3]{CompletenessV2},
generated cover edges~\cite[\S 4]{CompletenessV2}, support-closed
mosaics~\cite[\S 6]{CompletenessV2}, and reachability
discharge~\cite[\S 11]{CompletenessV2}.

\smallskip
\noindent\emph{Effective enumeration.} By v1 finite-index lemma
(Thm~1.10), the set $\Sigma_d$ of rank-$d$ states is finite; hence
finitely many candidate quotients exist up to isomorphism. The v2
validity conditions (v1 Def~1.10 with Def~1.7 and Def~1.5
replaced as above) are decidable on a finite quotient: each of
(C1)--(C5) is a finite local check; reachability discharge is
finite-graph reachability; the witness menu's
$\Need_R$-filtering for $\DR/\PO$ existentials is a finite check
on the disjunctive network alphabet.

\smallskip
Therefore: enumerate finite quotients up to isomorphism, check
v2 validity, accept iff a valid quotient with $C_0 \in
\lambda(\sigma_\ast)$ is found. By soundness, acceptance implies
$C_0$ is satisfiable. By completeness, satisfiability of $C_0$
implies a valid quotient exists. The procedure is decidable.
\qedhere
\end{proof}

\subsection{What \emph{is not} claimed}

\begin{itemize}
\item No complexity bound. The quotient enumeration is
  non-elementary in $|C_0|$ on its face. An $\textsc{ExpTime}$
  upper bound (matching the $\mathcal{ALCI}$ lower bound) would
  require a tighter analysis of $|\Sigma_d|$, which is open.
  v1 simply asserted finiteness of $\Sigma_d$ without a
  quantitative bound; v2 does the same.
\item No claim about $\ALCIRCC{8}$. Both v1 and v2 only treat the
  RCC5 case. Extension to RCC8 (which also has the patchwork
  property of Renz \& Nebel~\cite{RenzNebel}) is expected to
  carry through by symmetric reasoning, but has not been written
  out.
\item No formal equivalence with the cover-tree tableau
  implementation~\cite{ClaudeImpl}. The implementation is an
  algorithmic realisation that is empirically validated on
  911 + 400 cross-validated test concepts and 12.2M enumerated
  models, but the formal equivalence between its CT1--CT5 checks
  and v2 quotient validity is a separate open claim.
\end{itemize}

\section{Comparison with v1 (line-level)}
\label{sec:comparison}

For ease of cross-reference with v1, this section lists the
specific edits, by v1 line number, that v2 induces.

\begin{center}
\renewcommand{\arraystretch}{1.25}
\begin{tabular}{p{0.27\linewidth}p{0.32\linewidth}p{0.32\linewidth}}
\toprule
v1 location & v1 reading & v2 reinterpretation \\
\midrule
Def 1.5 (line 161) & ``rooted tree~$T$ with ambient root~$\varrho_T$'' &
no ambient root; $\mathsf{par}$ may self-loop at root
(Def~\ref{def:gen-split-tree}) \\
Def 1.5 (line 167) & ``tree edges are immediate $\PPI$ edges'' &
generated witness-cover edges relative to $A = \orig(T)$
(Remark~\ref{rem:generated-covers}) \\
Def 1.7 (line 367) & $\Wit_k$ records target color for each
existential, including $\PP/\PPI$ &
$\Wit_k$ retains $\DR/\PO$ slots only;
$\PP/\PPI$ discharged by reachability
(Def~\ref{def:reach-discharge}) \\
Thm 1.20 (line 681--691) & decidability proof sketch with
ambient-root choice and ``least support-closed tables'' assertion &
proof replaced by Theorem~\ref{thm:decidability-v2}; model-to-quotient
delegated to~\cite[Thm 12.1]{CompletenessV2} \\
Remark following Thm 1.20 (line 693--695) & ``the only compressed
part is the extraction of support-closed tables'' &
extraction is now explicit~\cite[\S 6, \S 9]{CompletenessV2} \\
\bottomrule
\end{tabular}
\end{center}

All other line-by-line content of v1 stands.

\section{Conclusion}

GPT-5.4's v1 sibling-interface paper~\cite{SiblingV1} contains
a soundness chain (Thms 1.17--1.19) for the split-forest
decidability programme that GPT-5.5's review~\cite{GPT55Review}
explicitly listed as salvageable, alongside two structural
defects in adjacent material: an ambient-root assumption in
Def~1.5 that fails on $C_\uparrow$, and a Thm~1.20 proof sketch
that asserts but does not construct support-closed tables and
conflates witnesses across distinct occurrences when discharging
the C5 axiom. v2 replaces Def~1.5 with a generated witness-cover
presentation (parent self-loops permitted), replaces the
witness-menu treatment of $\PP/\PPI$ existentials by
reachability discharge on the finite parent graph, and replaces
the Thm~1.20 sketch with a citation to v2 completeness
extraction~\cite{CompletenessV2}. The soundness chain of
Thms~1.17--1.19 is preserved verbatim. Combined with v2
completeness extraction, this yields a decidability theorem for
$\ALCIRCC{5}$ concept satisfiability under strong-$\EQ$
semantics, without B\"uchi machinery or
$\omega$-acceptance.

\begin{thebibliography}{9}

\bibitem{SiblingV1}
GPT-5.4 (OpenAI).
\newblock Finite sibling-interface descriptors and a proposed
  completion for $\mathcal{ALCI}_{RCC5}$.
\newblock Proof-style note, April 2026. Prompted by Michael Wessel.
\newblock Filename in repository:
  \texttt{papers/trees/sibling\_interface\_descriptors\_ALCIRCC5\_completed\_eqsync\_canonical\_needpatched.tex}.
\newblock Superseded in part by the present v2.

\bibitem{GPT55Review}
GPT-5.5 (OpenAI).
\newblock Review of the $\mathcal{ALCI}_{RCC5}$ decidability
  manuscripts.
\newblock Review note, May 2026. Prompted by Michael Wessel.
\newblock Filename in repository:
  \texttt{papers/gpt5.5/gpt5.5-original-review.pdf}.

\bibitem{GPT55Generated}
GPT-5.5 (OpenAI).
\newblock A coherent generated split-forest decidability proof for
  $\mathcal{ALCI}_{RCC5}$.
\newblock Proof manuscript, May 2026. Prompted by Michael Wessel.
\newblock Filename in repository:
  \texttt{papers/gpt5.5/ALCIRCC5\_coherent\_generated\_split\_forest\_decidability.pdf}.

\bibitem{GPT55Survival}
GPT-5.5 (OpenAI).
\newblock What survives of the split-forest idea.
\newblock Companion note to the repaired decidability proof, May 2026.
\newblock Filename in repository:
  \texttt{papers/gpt5.5/split\_forest\_survival\_note.pdf}.

\bibitem{CompletenessV1}
Claude (Anthropic).
\newblock Completeness extraction for $\mathcal{ALCI}_{RCC5}$ (v1).
\newblock Proof-style note, April 2026. Prompted by Michael Wessel.
\newblock Superseded by v2~\cite{CompletenessV2}.

\bibitem{CompletenessV2}
Claude (Anthropic, Opus 4.7).
\newblock Completeness extraction for $\mathcal{ALCI}_{RCC5}$ (v2,
  repaired).
\newblock Proof-style note, May 2026. Prompted by Michael Wessel.
\newblock Filename in repository:
  \texttt{papers/completeness\_extraction\_ALCIRCC5\_v2.tex}.

\bibitem{CoverTreeTableauPatched}
GPT-5.4 (OpenAI).
\newblock Cover-tree tableau for $\mathcal{ALCI}_{RCC5}$
  (need-all patched).
\newblock Proof-style note, April 2026.
\newblock Filename in repository:
  \texttt{papers/trees/alcircc5\_cover\_tree\_tableau\_needall\_patched.tex}.

\bibitem{ClaudeImpl}
Claude (Anthropic).
\newblock Cover-tree tableau implementation for $\mathcal{ALCI}_{RCC5}$.
\newblock Implementation paper, April 2026.

\bibitem{RenzNebel}
J. Renz and B. Nebel.
\newblock On the complexity of qualitative spatial reasoning: a maximal
  tractable fragment of the region connection calculus.
\newblock \emph{Artificial Intelligence}, 108(1--2):69--123, 1999.

\bibitem{Wessel2003}
M. Wessel.
\newblock Qualitative spatial reasoning with the
  $\mathcal{ALCI}_\mathit{RCC}$-family --- first results and
  unanswered questions.
\newblock Technical Report FBI-HH-M-324/03, University of Hamburg,
  2002/2003.

\end{thebibliography}

\end{document}
